\documentclass[11pt]{article}
\usepackage{style/blueprintdoc}
\geometry{margin=14mm,top=15mm,bottom=13mm}
\setdoccode{NNV-JS}\setdoctitle{Judging scoresheet — responsible AI projects}
\begin{document}
\bptitleblock{Judging scoresheet}{One sheet per team · suggested criteria, to be aligned with the panel's rubric}{DOCUMENT NNV-JS \quad|\quad REV A \quad|\quad Team: \rule{4.5cm}{0.4pt} \quad Judge: \rule{3cm}{0.4pt}}
\footnotesize
\begin{tabularx}{\textwidth}{@{}>{\bfseries}p{3.1cm}Xc@{}}
\toprule
criterion & what a 5 looks like (a 1 is the opposite) & score 1–5\\
\midrule
Problem and people & A real need, named end users or stakeholders who were actually consulted; the team can say who is affected when the system is wrong. & \rule{0pt}{5mm}\_\_\_\\
Principle made concrete & The responsible-AI principle(s) claimed (fairness, transparency, privacy, safety, governance, sustainability) show up in the \emph{product}: a design choice, a measurement, a control — not only in the slides. & \rule{0pt}{5mm}\_\_\_\\
Evidence, with its scope & Every number comes with the inputs it was measured on. ``96\,\% accurate'' says \emph{on what}. No claim is wider than its evidence. & \rule{0pt}{5mm}\_\_\_\\
Assurance argument & At least one requirement is stated as a property over a region of inputs (``for every \ldots\ in \ldots, the system must \ldots''), with how it was — or could be — checked, and what a counterexample would look like. & \rule{0pt}{5mm}\_\_\_\\
Honesty about limits & What is out of scope, which assumptions are load-bearing, what fails first and what happens then (human oversight, monitoring, fallback). & \rule{0pt}{5mm}\_\_\_\\
Accountability & Who chooses the thresholds, who signs off, and how an affected person finds out and contests a decision. & \rule{0pt}{5mm}\_\_\_\\
Progress and craft & A visible MVP → Alpha → Beta → final trajectory; something that works end to end; feedback from mentors and users was used. & \rule{0pt}{5mm}\_\_\_\\
Communication & Clear to a non-specialist in the room; the team distinguishes what they know from what they hope. & \rule{0pt}{5mm}\_\_\_\\
\midrule
\multicolumn{2}{r}{\textbf{total (of 40)}} & \rule{0pt}{5mm}\_\_\_\\
\bottomrule
\end{tabularx}
\vspace{2mm}
\begin{notebox}\footnotesize
\textbf{Tie-breakers.} Prefer the team that (1) can name the counterexample that would change their mind, (2) shows a decision sheet or worksheet with the blanks filled in — including \emph{out of scope}, and (3) turned a stakeholder's objection into a change in the product.
\end{notebox}
\vspace{1mm}
\begin{tcolorbox}[enhanced,colback=white,colframe=bpRule,boxrule=0.5pt,arc=0pt,height=4.6cm,left=2mm,right=2mm,top=1mm,bottom=1mm,fontupper=\footnotesize]
\textbf{Notes for feedback} (one strength, one risk, one question to take home):
\end{tcolorbox}
\end{document}
